\chapter{Reproducibility and Supplementary Material}
\label{chap:appendix_a}

\section{Code and Data Availability}
\label{sec:availability}

The complete pipeline --- data loading, feature engineering, the four models,
training, evaluation and all plotting --- is released as open source at
\begin{center}
\url{https://github.com/ramc77/fyp-jet-tagging}
\end{center}
The dataset is the public Top-Quark Tagging Reference
Dataset~\citep{kasieczka2019dataset}, available from Zenodo
(\url{https://zenodo.org/records/2603256}).

\section{How to Reproduce the Results}
\label{sec:reproduce}

On a machine with Python~3.11 the headline results and study are reproduced
with:
\begin{verbatim}
python3 -m venv particle_venv
source particle_venv/bin/activate
pip install -r requirements.txt
python3 download_data.py --yes        # fetch the dataset
python3 run_full_pipeline.py          # train + evaluate all four models
python3 run_study.py                  # data-efficiency + calibration sweep
\end{verbatim}
A Google Colaboratory notebook (\texttt{notebooks/FYP\_Colab.ipynb}) runs the
same pipeline on a free T4 GPU and downloads the results. Every training stage
checkpoints after each epoch and resumes automatically, so a run may be
interrupted and continued.

\section{Training Configuration}
\label{sec:config}

\begin{table}[htbp]
\centering
\caption{Principal hyper-parameters of the four models.}
\label{tab:hyperparams}
\ra{1.2}
\begin{tabular}{@{}lll@{}}
\toprule
\textbf{Model} & \textbf{Key settings} & \textbf{Parameters} \\
\midrule
BDT (XGBoost)        & 500 rounds, depth 7, lr 0.1   & $\sim$160 trees \\
CNN (jet images)     & 3 ResNet blocks, 64--256 ch.  & $1.26$\,M \\
ParticleNet          & 3 EdgeConv blocks, $k=16$     & $0.37$\,M \\
Particle Transformer & 3 blocks, 4 heads, dim 128    & $0.43$\,M \\
\bottomrule
\end{tabular}
\end{table}

All models minimise binary cross-entropy with the Adam/AdamW optimiser,
gradient clipping, a plateau learning-rate scheduler and early stopping on
the validation AUC, as detailed in Section~\ref{sec:training}.

\section{Training History}
\label{sec:history}

Figure~\ref{fig:history} shows the training and validation loss and AUC
curves for the three neural models, confirming smooth convergence without
significant over-fitting within the epoch budget.

\begin{figure}[htbp]
\centering
\includegraphics[width=\textwidth]{training_history.pdf}
\caption{Training/validation loss and validation AUC versus epoch for the
CNN, ParticleNet and Particle Transformer.}
\label{fig:history}
\end{figure}

\section{Transformer Attention Maps}
\label{sec:attnmaps}

Figure~\ref{fig:attnmaps} shows representative head-averaged self-attention
maps for top jets across the transformer layers, complementing the
class-token analysis in Section~\ref{sec:res_interpret}.

\begin{figure}[htbp]
\centering
\includegraphics[width=0.9\textwidth]{attention_maps.pdf}
\caption{Head-averaged self-attention maps of the Particle Transformer for
sample top jets across the encoder layers. The padded (non-particle)
positions appear as uniform blocks.}
\label{fig:attnmaps}
\end{figure}
